%!TEX root = ../graduate-work.tex
\section{Experiments}

\subsection{Methodology and Experimental Plan}

Eight experiments were conducted to validate the implemented methods, grouped
by technique and the property under test. Experiments~1--3 concern Tensor
Parallelism: memory consumption when performing CROWN verification of a model
that does not fit on a single GPU; correctness and soundness of bounds in the
presence of multiple sharded zones; automatic sharding of ONNX models from
VNN-COMP~2022. Experiments~4--8 concern FSDP: bitwise identity of bounds;
savings in baseline and peak memory as a function of model width and depth;
fair comparison of FSDP$=$2 vs.\ single-GPU in $\beta$-CROWN+BaB under an
identical workload; verification of ViT after resolving JIT incompatibilities;
and correctness of BoundConv sharding on convolutional models from VNN-COMP'24.

Experiments~1--5 were conducted on Server~A: 2$\times$~NVIDIA~A40 (48~GB VRAM,
Ampere), inter-GPU connection via NVLink. Experiments~6--8 were conducted on
Server~B: 2$\times$~NVIDIA~A40 (48~GB), connection via PCIe without NVLink;
the flag \texttt{NCCL\_P2P\_DISABLE=1} was set to prevent AllGather hangs.
The collective-operation backend was NCCL; framework was PyTorch~2.x; all
distributed runs were launched via \texttt{torchrun --nproc\_per\_node=2}.

\subsection{Experiment 1: GPU Memory Consumption}

\subsubsection{Model Architecture}

The target of verification was a fully connected network with a single hidden
layer:
$$f(\mathbf{x}) = W_2 \cdot \mathrm{ReLU}(W_1 \mathbf{x} + \mathbf{b}_1) + \mathbf{b}_2,$$
where $W_1 \in \mathbb{R}^{H \times D}$, $W_2 \in \mathbb{R}^{1 \times H}$,
$D = 4096$, $H = 262144$. The matrix $W_1$ of size $262144 \times 4096$
($\approx\!4$~GB in float32) dominates memory consumption both during the
forward pass and during CROWN bound computation.

\subsubsection{Tensor Parallelism Implementation}

For the TP variant, the model was partitioned following the Megatron-LM
scheme~\cite{MegatronLM2020}. $W_1$ is column-sharded: GPU~$r$ holds
$W_1^{(r)} \in \mathbb{R}^{(H/P) \times D}$. $W_2$ is row-sharded: GPU~$r$
holds $W_2^{(r)} \in \mathbb{R}^{1 \times (H/P)}$, computes a partial result,
and \texttt{AllReduce} reconstructs the full output. For integration with
\texttt{auto\_LiRPA}~\cite{AutoLiRPARepo}, operators \texttt{BoundLinearTP\_Col}
and \texttt{BoundLinearTP\_Row} were implemented, adding \texttt{AllReduce} to
the bound backward pass.

\subsubsection{Methodology}

Verification was performed using CROWN for $L_\infty$-robustness with
$\varepsilon = 0.01$ on a batch of $N = 2048$ random inputs of dimension
$D = 4096$. Peak GPU memory was measured via
\texttt{torch.cuda.max\_memory\_allocated}.

\subsubsection{Results}

\begin{table}[htbp]
\centering
\caption{Peak GPU memory during CROWN verification
($D{=}4096$, $H{=}262144$, $N{=}2048$, $\varepsilon{=}0.01$)}
\label{tab:tp-memory}
\begin{tabular}{lccc}
\toprule
Mode & GPU & \texttt{max\_allocated}, MB & \texttt{max\_reserved}, MB \\
\midrule
Single GPU & 1$\times$A40 & 26\,826 & 30\,874 \\
TP = 2     & GPU 0        & 13\,513 & 15\,512 \\
TP = 2     & GPU 1        & 13\,513 & 15\,512 \\
\bottomrule
\end{tabular}
\end{table}

The memory reduction factor is:
$$k = \frac{26\,826}{13\,513} \approx 1.985,$$
close to the theoretical limit of $P = 2$. The minor deviation is explained
by non-shardable structures: the input tensor $\mathbf{x} \in \mathbb{R}^{N \times D}$,
NCCL buffers, and computational graph metadata are replicated on every GPU.

\subsection{Experiment 2: Correctness and Soundness of Bounds}

\subsubsection{Methodology}

For each test model, bounds were computed two ways: standard single-GPU CROWN
(full weight matrix, no \texttt{AllReduce}) and CROWN with TP=2 (sharded
matrices with \texttt{AllReduce}). Three properties were verified: maximum
deviation $|lb_{\text{ref}} - lb_{\text{tp}}|$ and
$|ub_{\text{ref}} - ub_{\text{tp}}|$; soundness
($lb_{\text{tp}} \leq lb_{\text{ref}}$ and $ub_{\text{tp}} \geq ub_{\text{ref}}$
for all outputs); and cross-rank consistency (both GPUs produce identical bounds).
Tests covered PyTorch MLP models of varying depth and ONNX models from
VNN-COMP~2022 (MNIST-FC)~\cite{arXiv2212VNNCOMP2022};
$\varepsilon = 0.02$, $L_\infty$-robustness.

\subsubsection{Results}

\begin{table}[htbp]
\centering
\caption{Correctness of TP=2 CROWN bounds relative to the single-GPU reference}
\label{tab:tp-correctness}
\begin{tabular}{lccccc}
\toprule
Model & Col-Row pairs & $|lb_{\text{diff}}|$ & $|ub_{\text{diff}}|$
    & Cross-rank & Sound \\
\midrule
PyTorch 256$\times$2 & 1 & $2.98 \cdot 10^{-8}$ & $2.98 \cdot 10^{-8}$ & 0 & Yes \\
PyTorch 256$\times$4 & 2 & $2.51 \cdot 10^{0}$  & $2.49 \cdot 10^{0}$  & 0 & Yes \\
ONNX 256$\times$2    & 1 & $5.96 \cdot 10^{-8}$ & $5.96 \cdot 10^{-8}$ & 0 & Yes \\
ONNX 256$\times$4    & 2 & $5.64 \cdot 10^{0}$  & $3.98 \cdot 10^{0}$  & 0 & Yes \\
ONNX 256$\times$6    & 3 & $7.92 \cdot 10^{2}$  & $7.50 \cdot 10^{2}$  & 0 & Yes \\
\bottomrule
\end{tabular}
\end{table}

\subsubsection{Analysis}

Models with one Column--Row pair (256$\times$2) yield bounds that agree with
the reference up to machine epsilon ($\sim\!10^{-8}$), confirming arithmetic
correctness: \texttt{AllReduce} in CROWN backward exactly reproduces the result
of a full matrix multiplication.

With multiple sharded zones, bounds become wider than the reference (soundness
is preserved but tightness degrades). The loss grows with the number of zones
and is caused not by arithmetic errors but by the method substitution:
intermediate bounds within zones are computed via IBP, since the CROWN backward
cannot cross zone boundaries with partial weight matrices.

Cross-rank deviation is zero in all tests, confirming deterministic computation.

\subsection{Experiment 3: Automatic Sharding of VNN-COMP Models}

This experiment tested \texttt{tp\_shard\_bounded\_module} on models from
VNN-COMP~2022~\cite{arXiv2212VNNCOMP2022}. MNIST-FC models were loaded in ONNX
format, converted to PyTorch via \texttt{onnx2pytorch}, wrapped in
\texttt{BoundedModule}, and automatically sharded---without any manual
intervention. The function traverses the \texttt{BoundedModule} graph in
topological order, locates alternating pairs of \texttt{BoundLinear} nodes, and
replaces them with \texttt{BoundLinearTP\_Col/Row} with sharded weights; it
then re-executes the forward pass to update graph metadata and marks sharded
zones for correct intermediate bound computation.

All three VNN-COMP models (256$\times$2, 256$\times$4, 256$\times$6) were
sharded successfully and passed the soundness check
(Table~\ref{tab:tp-correctness}), confirming the applicability of the approach
to arbitrary fully connected ONNX models.

\subsection{Experiment 4: Correctness of FSDP Bounds}

\subsubsection{Methodology}

The methodology mirrors Experiment~2: bounds were computed single-GPU and
FSDP$=$2 with verification of deviations, cross-rank consistency, and soundness.
Tests covered MLPs ($h = 256$, depth 2, 4, 6 layers) under IBP and CROWN, plus
VNN-COMP ONNX models (256$\times$2, 256$\times$4) under CROWN---eight tests in total.

\subsubsection{Results}

\begin{table}[htbp]
\centering
\caption{Correctness of FSDP=2 bounds relative to the single-GPU reference.
$\Delta lb$, $\Delta ub$---maximum absolute deviation from reference bounds.}
\label{tab:fsdp-correctness}
\begin{tabular}{llcccc}
\toprule
Model & Method & $|\Delta lb|$ & $|\Delta ub|$
    & Cross-rank & Sound \\
\midrule
MLP 256$\times$2 & IBP   & $0.00$ & $0.00$ & $0.00$ & Yes \\
MLP 256$\times$2 & CROWN & $0.00$ & $0.00$ & $0.00$ & Yes \\
MLP 256$\times$4 & IBP   & $0.00$ & $0.00$ & $0.00$ & Yes \\
MLP 256$\times$4 & CROWN & $0.00$ & $0.00$ & $0.00$ & Yes \\
MLP 256$\times$6 & IBP   & $0.00$ & $0.00$ & $0.00$ & Yes \\
MLP 256$\times$6 & CROWN & $0.00$ & $0.00$ & $0.00$ & Yes \\
ONNX 256$\times$2 & CROWN & $0.00$ & $0.00$ & $0.00$ & Yes \\
ONNX 256$\times$4 & CROWN & $0.00$ & $0.00$ & $0.00$ & Yes \\
\bottomrule
\end{tabular}
\end{table}

\subsubsection{Analysis}

FSDP yields \textit{bitwise-identical} bounds for all tests: the deviation is
exactly $0.00$---a true zero, not machine epsilon. This is a fundamental
distinction from TP (Table~\ref{tab:tp-correctness}), where models with multiple
zones lose tightness by orders of magnitude.

The reason is that each layer under FSDP receives the full weight matrix via
\texttt{AllGather} before computation begins, so IBP and CROWN are
arithmetically identical to the single-GPU case. Sharding affects only weight
storage, not the operations themselves.

\subsection{Experiment 5: GPU Memory Consumption with FSDP}

\subsubsection{Methodology}

Two metrics were measured. \textit{Baseline memory}---the allocation after
initialisation and sharding, before \texttt{compute\_bounds} is called;
reflects the cost of weight storage. \textit{Peak memory}---
\texttt{max\_memory\_allocated} during \texttt{compute\_bounds}; includes
weights, CROWN $A$-matrices, intermediate tensors, and NCCL buffers.

Tests covered MLPs with width $h \in \{256, 1024, 4096, 8192\}$ at depth~4,
and depth $d \in \{2, 4, 6, 8\}$ at width~4096. Verification used CROWN,
$\varepsilon = 0.02$, input dimension~784 (MNIST). For FSDP$=$2, peak memory
was taken as the maximum over both ranks.

\subsubsection{Results}

\begin{table}[htbp]
\centering
\caption{Baseline memory (weight storage) with FSDP=2 vs.\ single-GPU}
\label{tab:fsdp-baseline}
\begin{tabular}{lrrr}
\toprule
Model & Single GPU, MB & FSDP, MB & Savings, \% \\
\midrule
MLP $h{=}256$, $d{=}4$   &     1.5 &    0.8 & 50 \\
MLP $h{=}1024$, $d{=}4$  &    15.1 &    7.6 & 50 \\
MLP $h{=}4096$, $d{=}4$  &   204.5 &  102.3 & 50 \\
MLP $h{=}8192$, $d{=}4$  &   792.9 &  396.5 & 50 \\
MLP $h{=}4096$, $d{=}8$  &   460.5 &  230.3 & 50 \\
\bottomrule
\end{tabular}
\end{table}

\begin{table}[htbp]
\centering
\caption{Peak memory (during \texttt{compute\_bounds}) with FSDP=2 vs.\ single-GPU}
\label{tab:fsdp-peak}
\begin{tabular}{lrrr}
\toprule
Model & Single GPU, MB & FSDP, MB & Savings, \% \\
\midrule
MLP $h{=}4096$, $d{=}4$  &   3\,800 &  2\,527 & 34 \\
MLP $h{=}8192$, $d{=}4$  &  14\,759 &  9\,782 & 34 \\
MLP $h{=}4096$, $d{=}8$  &  16\,726 & 10\,232 & 39 \\
\bottomrule
\end{tabular}
\end{table}

\subsubsection{Analysis}

Baseline memory is reduced by exactly 50\% for every configuration: each rank
stores exactly half the rows of every weight tensor, matching the theoretical
value $1/P = 50\%$ (Section~\ref{sec:fsdp-memory-theory}).

Peak memory decreases by only 17--39\%, because FSDP shards only weights while
CROWN $A$-matrices always have full dimension $O(h \times h)$ on each GPU.
At $d = 4$ layers, $A$-matrices are comparable in size to weights, so sharding
weights alone yields ${\sim}34\%$; at $d = 8$, the weight fraction is higher
and savings reach 39\%. For shallow models ($d = 2$) peak savings are minimal
since $A$-matrices dominate.

For comparison, TP (Exp.~1) achieves ${\sim}2\times$ peak memory reduction by
sharding both weights and $A$-matrices, at the cost of bound degradation.
FSDP offers moderate savings with perfect bound accuracy---critical for complete
verification.

\subsection{Experiment 6: FSDP in Complete Verification (BaB)}
\label{sec:exp6}

\subsubsection{Motivation and Methodology Correction}

Complete verification ($\beta$-CROWN~\cite{betaCROWN2021} + BaB) differs
substantially from incomplete verification: it iteratively splits the input
domain into sub-problems, optimises $\alpha$- and $\beta$-parameters for each
sub-problem, and stores the branching-tree history in GPU memory. The primary
memory consumers here are \textit{alpha tensors}---per-neuron ReLU approximation
slopes $\alpha_{i,k} \in [0, 1]$, with dimensionality
$(\text{layers}) \times (\text{specifications}) \times (\text{batch}) \times
(\text{neurons per layer})$.

During the development of a fair comparison, a systematic artefact was discovered:
the single-GPU mode used \texttt{auto\_enlarge\_batch\_size} (automatically
filling the GPU) and \texttt{pruning\_in\_iteration} (early round termination),
whereas FSDP mode disabled these optimisations for AllGather synchronisation.
The two modes were thus processing \textit{different workloads}---the apparent
``27$\times$'' advantage of FSDP in early runs was spurious.

To eliminate the artefact, the option \texttt{bab.force\_synchronous}
(\texttt{--force\_synchronous}) was implemented, which forces single-GPU mode to
follow the same BaB code path as FSDP: fixed batch size, no
\texttt{auto\_enlarge\_batch\_size}, no \texttt{early\_stop}, no
\texttt{pruning\_in\_iteration}.

\subsubsection{Methodology}

The model \texttt{mnist-net\_256x6} (6 layers of width 256,
${\sim}1.5$~MB weights) from VNN-COMP~2022 was used with $\varepsilon = 0.05$.
Parameters: exactly 10 BaB rounds, batch sizes of 512 and 4096, identical for
single-GPU and FSDP$=$2.

\subsubsection{Results}

\begin{table}[htbp]
\centering
\caption{Fair comparison of FSDP$=$2 and single-GPU in BaB on
\texttt{mnist-net\_256x6}, $\varepsilon{=}0.05$, exactly 10 rounds.
Workload is identical.}
\label{tab:fsdp-bab}
\begin{tabular}{lcccc}
\toprule
Configuration & Batch & Domains & Rounds & Peak GPU, MB \\
\midrule
Single GPU  &   512 &   6\,236 & 10 &     211.7 \\
FSDP$=$2    &   512 &   6\,236 & 10 & 215.7 / rank \\
Single GPU  & 4\,096 &  49\,888 & 10 &   1\,537.3 \\
FSDP$=$2    & 4\,096 &  49\,888 & 10 & 1\,541.4 / rank \\
\bottomrule
\end{tabular}
\end{table}

\subsubsection{Analysis}

\texttt{mnist-net\_256x6} has ${\sim}1.5$~MB weights while total peak
consumption is ${\sim}1.5$~GB: weights constitute less than $0.1\%$; the
remainder is alpha/beta tensors and sub-domain history. Sharding $1.5$~MB does
not yield measurable gain; \texttt{AllGather} adds ${\sim}4$~MB overhead.

FSDP does not compromise bound accuracy and does not alter the complete
verification workload, but it does not save memory on models where weights
are negligible compared to alpha tensors.

\subsection{Experiment 7: FSDP with ViT from VNN-COMP'23}
\label{sec:exp7}

\subsubsection{Model and JIT Compatibility}

The model \texttt{pgd\_2\_3\_16} (2 Transformer blocks, 3 heads,
${\sim}75$K parameters, ${\sim}0.3$~MB weights) from VNN-COMP 2023 was used.
Unlike MLPs, the Transformer employs dynamic reshape operations that are
incompatible with the JIT tracing used by \texttt{auto\_LiRPA}.

Three classes of JIT incompatibilities were resolved. In \texttt{BoundReshape},
the batch dimension is ``baked in'' as a constant during JIT tracing; the fix
replaces the first axis with the current batch size whenever
$\prod(\text{new\_shape}) = x.\text{numel}()$---the check prevents incorrect
replacement for reshapes of the form $[B, N, D] \to [B \cdot N, D]$.
In \texttt{BoundConcat}, the JIT fixed the number of concatenated tensors.
In \texttt{BoundConstantOfShape}, the JIT-baked shape is replaced with the
current shape.

\subsubsection{Results}

\begin{table}[htbp]
\centering
\caption{Fair comparison of FSDP$=$2 and single-GPU on ViT
\texttt{pgd\_2\_3\_16}, VNN-COMP'23}
\label{tab:fsdp-vit}
\begin{tabular}{lccc}
\toprule
Configuration & Domains & Rounds & Peak GPU, MB \\
\midrule
Single GPU & 78 & 10 & 1\,020.1 \\
FSDP$=$2   & 78 & 10 & 1\,020.1 / rank \\
\bottomrule
\end{tabular}
\end{table}

The same pattern as in Experiment~6: weights ($0.3$~MB) are negligible compared
to attention $A$-matrices (${\sim}1$~GB). The key result is the \textit{correct
execution} of ViT verification in $\alpha,\beta$-CROWN with FSDP: bounds are
bitwise identical to single-GPU, and the workload is reproducible.

\subsection{Experiment 8: Extending FSDP to Convolutional Layers}
\label{sec:exp8}

\subsubsection{Implementation}

The previous experiments covered only fully connected layers. To support
vision models, FSDP was extended to \texttt{BoundConv}. Convolutional weight
tensors have shape $[C_\text{out}, C_\text{in}, k_H, k_W]$; sharding is
performed along dimension~0 (output channels $C_\text{out}$)---the only axis
along which the layer is linear in the output and along which \texttt{AllGather}
correctly reconstructs the full tensor.

Changes touched three files: \texttt{fsdp\_utils.py} (added \texttt{BoundConv}
to the list of shardable types and prevented double-sharding);
\texttt{backward\_bound.py} (added AllGather/free for \texttt{BoundConv} in
CROWN backward); and \texttt{interval\_bound.py} (added free hooks after IBP
propagation).

\subsubsection{Methodology}

ResNets from VNN-COMP~2024 (CIFAR-100) were used:
\texttt{CIFAR100\_resnet\_medium} and \texttt{CIFAR100\_resnet\_large}.
Fair conditions: \texttt{force\_synchronous}, batch$=$64, max\_iterations$=$5.

\subsubsection{Results}

\begin{table}[htbp]
\centering
\caption{Fair comparison of FSDP$=$2 and single-GPU on CIFAR-100 ResNet
(VNN-COMP'24), BoundConv sharding}
\label{tab:fsdp-conv}
\begin{tabular}{llcccc}
\toprule
Model & Config & Domains & Rounds & Peak, MB & $|\Delta lb|$ \\
\midrule
ResNet-medium & Single GPU & 130 & 5 &   960 & --- \\
ResNet-medium & FSDP$=$2   & 130 & 5 &   930 / rank & 0.00 \\
ResNet-large  & Single GPU & \textbf{unsat}\textsuperscript{$\dagger$} & 0 & 4\,610.6 & --- \\
ResNet-large  & FSDP$=$2   & \textbf{unsat}\textsuperscript{$\dagger$} & 0 & 4\,641.5 / rank & 0.00 \\
\bottomrule
\end{tabular}
\end{table}
ResNet-medium: budget exhausted (max\_iterations$=$5), result \emph{unknown}.
\textsuperscript{$\dagger$}ResNet-large: \textbf{property proved}---$\alpha$-CROWN
computed min-margin $\ge 0$ without branching (\emph{verified, unsat}).

\subsubsection{Analysis}

Bounds are bitwise identical for both models. For ResNet-large, a complete
verification result (unsat) was obtained: $\alpha$-CROWN proved the property
at initialisation without BaB. The agreement of results across both modes
confirms the correctness of BoundConv sharding.

ResNet-large contains ${\sim}95$~MB of convolutional weights, but alpha tensors
occupy over $4$~GB---they are not sharded, and each GPU stores the full set of
per-neuron slopes for the entire batch. The AllGather overhead (${\sim}30$~MB)
exceeds the savings from sharding $95$~MB of weights at $P=2$. In the
$\alpha$-CROWN+BaB regime, the memory bottleneck is not weights but alpha
tensors; sharding them is the next development step.
